\begin{helper}
Fix natural numbers \(p,m,t\), a prefix tuple
\[
  pref:\{0,\ldots,m-1\}\to \mathbb F_2^p\times\mathbb F_2^p,
\]
and a possible final first coordinate \(a\in\mathbb F_2^p\). Form the
first-coordinate list obtained by appending \(a\) to the prefix. The number of
vectors \(b\in\mathbb F_2^p\) such that this appended first-coordinate list has
rank \(t\) and
\[
  \langle a_{\ell+1},b\rangle=1
  \quad\text{for every zero-based }\ell<m+1
\]
is at most \(2^{p-t}\). In Lean the rank is computed by appending the pair
\((a,0)\), because \(\noderef{F2BadTupleRank}\) depends only on first
coordinates.
\end{helper}
\begin{proof}
Let \(ab_0\) be the length \(m+1\) tuple obtained from \(pref\) by appending
the pair \((a,0)\). The counted set of possible last coordinates \(b\) is empty
unless
\[
  \mathrm{F2BadTupleRank}(p,m+1,ab_0,m+1)=t.
\]
If this rank equality fails, the desired bound is immediate.

Assume now that the rank equality holds. The counted \(b\)'s inject into the
fiber
\[
  \{y\in\mathbb F_2^p:
      \langle (ab_0)_\ell{}_1,y\rangle=1
      \text{ for every zero-based }\ell<m+1\}.
\]
Indeed, a counted \(b\) already satisfies exactly these dot-product equations.
By \(\noderef{F2BadTuplePrefixFiberBound}\), applied to \(ab_0\) with prefix
length \(m+1\), this fiber has cardinality at most
\[
  2^{p-\mathrm{F2BadTupleRank}(p,m+1,ab_0,m+1)}.
\]
Using the assumed rank equality, this is \(2^{p-t}\). Hence the counted set of
possible last \(b\)'s has cardinality at most \(2^{p-t}\).
\end{proof}
